% Standalone problem document, generated from the adjacent README.md.
% Compile with XeLaTeX. Mathematical content is maintained in Markdown.
\documentclass[11pt,a4paper]{article}
\usepackage[fontsize=10.5pt]{scrextend}
\usepackage[margin=22mm,headheight=15pt,headsep=9mm,footskip=11mm]{geometry}
\usepackage{fontspec}
\setmainfont{texgyrepagella}[Extension=.otf,UprightFont=*-regular,BoldFont=*-bold,ItalicFont=*-italic,BoldItalicFont=*-bolditalic]
\setsansfont{texgyreheros}[Extension=.otf,UprightFont=*-regular,BoldFont=*-bold,ItalicFont=*-italic,BoldItalicFont=*-bolditalic]
\setmonofont{texgyrecursor}[Extension=.otf,UprightFont=*-regular,BoldFont=*-bold,ItalicFont=*-italic,BoldItalicFont=*-bolditalic,Scale=MatchLowercase]
\usepackage{amsmath,amssymb}
\usepackage{unicode-math}
\setmathfont{texgyrepagella-math.otf}
\usepackage{xcolor,microtype,enumitem,needspace,fancyhdr}
\usepackage{bookmark}
\definecolor{ink}{HTML}{17344B}
\definecolor{accent}{HTML}{197D80}
\definecolor{muted}{HTML}{596773}
\hypersetup{colorlinks=true,linkcolor=accent,urlcolor=accent,pdftitle={IV-01 - Two-vertex
certification of nonsingular sign
regularity},pdfauthor={Open Problems in Numerical Linear Algebra}}
\urlstyle{same}
\setlength{\parindent}{0pt}
\setlength{\parskip}{5pt plus 1pt minus 1pt}
\linespread{1.04}
\setlength{\emergencystretch}{3em}
\widowpenalty=10000
\clubpenalty=10000
\displaywidowpenalty=10000
\allowdisplaybreaks[1]
\setlist{nosep,leftmargin=1.5em}
\providecommand{\tightlist}{\setlength{\itemsep}{0pt}\setlength{\parskip}{0pt}}
\providecommand{\pandocbounded}[1]{#1}
\pagestyle{fancy}
\fancyhf{}
\fancyhead[L]{\sffamily\footnotesize\color{muted}OPEN PROBLEMS IN NUMERICAL LINEAR ALGEBRA}
\fancyhead[R]{\sffamily\bfseries\footnotesize\color{accent}IV-01}
\fancyfoot[L]{\parbox[b]{0.9\textwidth}{\sffamily\fontsize{8}{10}\selectfont\color{muted}Editorial ratings; a bounded search cannot certify that a problem remains unsolved.\\Literature check: 2026-09-10}}
\fancyfoot[R]{\sffamily\footnotesize\color{muted}\thepage}
\renewcommand{\headrulewidth}{0.3pt}
\renewcommand{\headrule}{\hbox to\headwidth{\color{accent}\leaders\hrule height \headrulewidth\hfill}}
\makeatletter
\renewcommand{\section}{\Needspace{5\baselineskip}\@startsection{section}{1}{0pt}{12pt plus 2pt minus 2pt}{4pt}{\sffamily\bfseries\large\color{ink}}}
\renewcommand{\subsection}{\Needspace{4\baselineskip}\@startsection{subsection}{2}{0pt}{9pt plus 1pt minus 1pt}{3pt}{\sffamily\bfseries\normalsize\color{ink}}}
\makeatother
\setcounter{secnumdepth}{0}
\begin{document}
{\sffamily\small\color{muted}Intervals and absolute value equations\par}
\vspace{3pt}
{\raggedright\hyphenpenalty=10000\exhyphenpenalty=10000\sffamily\bfseries\fontsize{21}{24}\selectfont\color{ink}Two-vertex
certification of nonsingular sign regularity\par}
\vspace{7pt}
{\sffamily\small\textbf{Difficulty:} challenging\quad\textbf{Importance:} interesting
to specialist\par}
{\sffamily\small\bfseries\color{accent}Status: Partially resolved\par}
\vspace{4pt}
{\color{accent}\hrule height 0.8pt}
\vspace{6pt}
\textbf{Area:} structured matrices and interval linear algebra

\textbf{Rating rationale:} Challenging reflects extending two-vertex
certification through zero minors in arbitrary dimension; specialist
impact is the efficient recognition of structured interval matrix
families.

\subsection{Problem statement}\label{problem-statement}

For \(\epsilon\in\{-1,1\}^n\), call a real \(n\times n\) matrix \(M\)
\textbf{nonsingular sign regular with signature \(\epsilon\)} if
\(\det M\ne0\) and every minor of order \(k\) has determinant \(d\)
satisfying \(\epsilon_k d\ge0\), for every \(1\le k\le n\).

Let \(n\ge5\) and let \(L,U\in\mathbb R^{n\times n}\) satisfy \(L\le U\)
entrywise. Define the two checkerboard vertex matrices

\[
A^-_{ij}=\begin{cases}L_{ij},&i+j\text{ even},\\U_{ij},&i+j\text{ odd},\end{cases}
\qquad
A^+_{ij}=\begin{cases}U_{ij},&i+j\text{ even},\\L_{ij},&i+j\text{ odd}.\end{cases}
\]

\subsubsection{Question}\label{question}

If \(A^-\) and \(A^+\) are nonsingular sign regular with the same
signature \(\epsilon\), must every real matrix \(M\) with
\(L\le M\le U\) be nonsingular sign regular with signature \(\epsilon\)?

The implication would certify a property of an entire interval of
matrices from only two of its vertices. Allowing zero minors is
essential to the remaining problem.

\subsection{References}\label{references}

Jürgen Garloff, Mohammad Adm, and Jihad Titi,
\href{https://www.reliable-computing.org/reliable-computing-22-pp-001-014.pdf}{\emph{A
Survey of Classes of Matrices Possessing the Interval Property and
Related Properties}}, Reliable Computing \textbf{22} (2016), 1--14,
Conjecture 3.1, p.~7. Adm and Garloff,
\href{https://doi.org/10.1007/s44146-026-00223-y}{\emph{Certification of
the Sign Regularity of Matrix Intervals}}, Acta Scientiarum
Mathematicarum, published January 17, 2026, §3, especially the question
following Theorem 3.2 and the discussion of admissible signatures.

\subsection{Earlier status evidence ---
2026-09-08}\label{earlier-status-evidence-2026-09-08}

The 2026 survey explicitly retains this question. It records positive
results for strictly sign regular matrices, totally nonnegative
matrices, tridiagonal matrices, and certain signatures covering all
orders at most four. Theorem 3.3 also covers intervals whose fixed
entries \(L_{ij}=U_{ij}\), if any, all have the same parity of \(i+j\);
fixed entries of both parities are allowed in the general question. It
also records counterexamples when nonsingularity is omitted. Searches
for sign-regular interval conjectures and subsequent 2026 results found
no resolution of the displayed implication.

\subsection{Audit update --- 2026-09-10}\label{audit-update-2026-09-10}

Rechecked the
\href{https://link.springer.com/article/10.1007/s44146-026-00223-y}{2026
survey}, §3, Theorems 3.2--3.3 and the intervening open question.
Tridiagonal and certain signature classes, as well as the specified
restrictions on fixed-entry parity, prove substantive subclasses even
for \(n\ge5\). The unrestricted nonsingular case remains explicitly open
there. Searches for subsequent sign-regular interval certification
results found no general resolution; the status now exposes these
partial results.
\end{document}
